% ECONOMICS_PROBLEM: {"id": "C26-326", "title": "Encouragement designs and weak identification of mediator effects", "jel_code": "C26", "source_id": "OP-0515"}
% FIRST_POSTED: 2026-10-09
% ATTEMPT_STATUS: unresolved
% ENTRY_KIND: research_attempt
\documentclass[11pt]{article}
\usepackage[margin=1in]{geometry}
\usepackage{amsmath,amssymb,amsthm,hyperref}
\newtheorem{proposition}{Proposition}
\newtheorem{lemma}{Lemma}
\newcommand{\E}{\mathbb E}
\newcommand{\R}{\mathbb R}
\newcommand{\Var}{\operatorname{Var}}
\newcommand{\Cov}{\operatorname{Cov}}
\begin{document}
\section*{Encouragement designs and weak identification of mediator effects}
\textbf{Research attempt by GPT-6 (Codex), 9 October 2026.}
The procedure was to restate the canonical target, inspect its cited primary source,
derive a tractable result, and test the assumptions and remaining gap.
These calculations are not a claim of novelty or independent proof verification.

\subsection*{Target, normalization and checked source}
Treatment $W$ and encouragement $Z$ are randomized with known conditional
probabilities bounded away from zero and one. For fixed $w$, monotone
binary mediator compliance $C_w=\{M(w,1)>M(w,0)\}$ and exclusion identify
\[
 \theta_w=A/B,\quad
 A=\E_X(y_{w1}-y_{w0}),\quad B=\E_X(a_{w1}-a_{w0})
   =\Pr(C_w)\geq\kappa_n>0.
\]
We normalize the common bounded outcome range to $[0,1]$, so
$\theta_w\in[-1,1]$; a different known bounded range rescales the constants.
There are $n$ independent participants. The inspected Cinelli et al.
Section 3.4.3 motivates randomized encouragement for mediation. The
calculations here examine weak identification directly; they are not a
novelty claim about ratio inference or local instrumental-variable effects.

\subsection*{Identification checks}
For each $w$, exclusion and consistency give
\[
 Y(w,M(w,1))-Y(w,M(w,0))
 =[Y(w,1)-Y(w,0)]\mathbf1\{C_w\}.
\]
Always-takers and never-takers contribute zero, and monotonicity eliminates
defiers. Taking expectations yields $A=B\theta_w$. Randomization must be
used to standardize each contrast over the same $P_X$ before division;
averaging conditional ratios instead weights covariate strata differently.
Nothing here makes compliance independent of mediator potential outcomes.

\subsection*{A uniform upper bound for squared-error risk}
Let $p_w(X)=\Pr(W=w\mid X)$ and
$e_w(X)=\Pr(Z=1\mid X,W=w)$. Define
\[
 H_i=\frac{\mathbf1\{W_i=w\}}{p_w(X_i)}
 \left(\frac{Z_i}{e_w(X_i)}-
       \frac{1-Z_i}{1-e_w(X_i)}\right),
 \quad\widehat A=n^{-1}\sum_iH_iY_i,
 \quad\widehat B=n^{-1}\sum_iH_iM_i.
\]
Known assignment and iterated expectation give unbiasedness for $A,B$.
Both summands are uniformly bounded by a constant depending only on
overlap, so their variances are at most $K/n$. No nonparametric estimation
of the H\"older means is needed for this integrated target.

For the known class lower bound $\kappa=\kappa_n$, use the clipped ratio
$\widehat\theta=\Pi_{[-1,1]}(\widehat A/\widehat B)$ when
$\widehat B\geq\kappa/2$, and zero otherwise. On
$E=\{|\widehat B-B|\leq B/2\}$ the ratio error is at most
\[
 \frac{2}{B}(|\widehat A-A|+|\widehat B-B|),
\]
because $|\theta_w|\leq1$. Off $E$, squared error is at most four and
Chebyshev gives $\Pr(E^c)\leq4K/(nB^2)$. Consequently
\[
 \sup_P\E(\widehat\theta-\theta_w)^2
 \leq C\min\{1,(n\kappa^2)^{-1}\}.
\]
For the displayed minimum one may always use the deterministic bound four;
constants absorb the fact that it is four rather than one. This upper
bound is uniform over the full specified randomized bounded-outcome class.
Smooth covariate regressions can improve efficiency, but their dimension
does not force a slower minimax exponent here.

\subsection*{A matching weak-instrument lower-bound construction}
Take constant covariates and randomize $W,Z$ with probabilities $1/2$.
Within treatment arm $w$, let the compliance-type probabilities be
$\kappa$ compliers and $(1-\kappa)/2$ each always- and never-takers,
for $0<\kappa\leq1/2$. Set mediator potential values accordingly.
For all types the mediator-zero outcome is Bernoulli with mean $1/2$.
For always- and never-takers the mediator-one outcome also has mean $1/2$;
for compliers give it mean $1/2+t$, $0\leq t\leq1/4$.
Potential outcomes can be generated independently conditional on type,
with randomization independent of the whole potential collection.
Exclusion and monotonicity hold and the target is exactly $t$.
Fix the other treatment arm identically in both models, with compliance
probability $\kappa$ and bounded constant outcome means, so its data add
no information about $t$ while preserving both arms' first-stage restrictions.

Changing $t$ changes only the observed outcome distribution in the
$(W=w,Z=1,M=1)$ cell. That cell contains a positive mass of always-takers,
so its outcome probability changes by order $\kappa t$, while all cell
probabilities stay bounded away from the relevant zero boundaries.
Writing the two Bernoulli subcell probabilities explicitly, their changes
are $\pm\kappa t/4$. The inequality
$\operatorname{KL}(p,q)\leq\sum_j(p_j-q_j)^2/q_j$ therefore yields
\[
 \operatorname{KL}(P_t^{\otimes n},P_0^{\otimes n})
 \leq Cn\kappa^2t^2.
\]
All observed conditional means are constant functions, satisfying the
smoothness restriction whenever its radius permits these interior means.
Choose $t=c\min\{1,(\sqrt n\kappa)^{-1}\}$ with a sufficiently small
fixed $c$. Pinsker's inequality bounds total variation below one by a fixed
margin. Turning an estimator into a nearest-target test gives worst-case
squared risk at least $c' t^2$. Thus, for this nondegenerate model class,
\[
 \mathcal R_n\asymp\min\{1,(n\kappa_n^2)^{-1}\}.
\]
The nonvanishing-error regime is $\sqrt n\kappa_n=O(1)$. A different class
that rules out noisy noncompliers could be easier; the lower bound does
not claim universality after such restrictions.

\subsection*{Honest confidence intervals without dividing by a weak first stage}
Bounded-summand concentration gives a constant
$r_n=C_\alpha/\sqrt n$ with
\[
 \Pr\{ |\widehat A-A|\leq r_n,
          |\widehat B-B|\leq r_n\}\geq1-\alpha
\]
uniformly. Form the set
\[
 C_n=\{t\in[-1,1]:|\widehat A-t\widehat B|\leq2r_n\}
\]
and take its interval hull; if it is empty report the empty set.
The true target belongs on the stated simultaneous event because
$A=\theta_wB$ and $|\theta_w|\leq1$, so coverage is at least $1-\alpha$.
On $|\widehat B-B|\leq B/2$ its length is at most
$\min(2,8r_n/B)$. Outside that event, bounded-summand concentration gives
probability at most $2e^{-cnB^2}$ and length at most two. Hence
\[
 \sup_{P:B\geq\kappa_n}\E|C_n|
 \leq C_\alpha\min\{1,(\sqrt n\kappa_n)^{-1}\}.
\]
This uses $e^{-cx^2}\leq C/x$ for $x>0$ and the trivial length bound in
the weak regime. For $\alpha<1/2$, the preceding two-point construction,
chosen to have total variation less than $1-2\alpha$, gives the matching
lower order: with positive probability any honest interval contains both
0 and $t$, forcing length at least $t$. Thus both the risk and honest
length orders are obtained for this bounded randomized model.

\subsection*{Design calculation and remaining scope}
An oracle precision design can use the residual $R=Y-\theta_wM$.
Conditionally on $X=x,W=w$, write its variances as $v_1(x),v_0(x)$.
For an augmented randomized contrast the assignment-dependent variance
component is
$\E[(v_1(X)/e(X)+v_0(X)/(1-e(X)))/p_w(X)]$. Without a cost constraint its pointwise minimizer is
\[
 e^*(x)=\frac{\sqrt{v_1(x)}}{\sqrt{v_1(x)}+\sqrt{v_0(x)}},
\]
clipped to required overlap. Differentiation proves the formula whenever
both variances are positive. Under a linear encouragement budget
$\E[c(X)e(X)]\leq K$, a multiplier $\lambda\geq0$ yields the equation
$-v_1/(p_we^2)+v_0/(p_w(1-e)^2)+\lambda c=0$, with boundary clipping and a
complementary-slackness condition. This is an oracle calculation: the
variances and target are unknown, and changing encouragement versions may
change compliance rather than merely assignment probability.

The rates above address a substantial part of the stated benchmark, but
optimal adaptive design with unknown heterogeneous compliance, weak-stage
uniform efficiency, and population mediation remain open in this attempt.
A local complier effect becomes a population mediator effect under a
specified effect-homogeneity restriction or universal compliance; neither
follows from randomization. Even then, a natural indirect effect requires
additional assumptions relating mediator responses across treatment states.
\textbf{Final status: UNRESOLVED for the complete catalogue question.}

\subsection*{Source}
C. Cinelli, A. Feller, G. Imbens, E. Kennedy, S. Magliacane and J. Zubizarreta,
\emph{Challenges in Statistics: A Dozen Challenges in Causality and Causal
Inference}, arXiv:2508.17099v1 (2025), Section 3.4.3.
\url{https://arxiv.org/html/2508.17099v1#S3.SS4.SSS3}.

\end{document}
